\documentclass[a4paper]{article}

\usepackage[polish]{babel}
\usepackage[T1]{fontenc}
\usepackage[utf8]{inputenc}
\usepackage{graphicx}
\usepackage{fancyhdr}
\usepackage{hyperref}
\pagestyle{fancy}

\lhead[]{}
\chead[Fundacja CoderDojo Polska]{Fundacja CoderDojo Polska}
\rhead[]{}
\cfoot[Treść udostępniana na licencji Creative Commons – Uznanie Autorstwa – Na Tych Samych Warunkach]{Treść udostępniana na licencji Creative Commons – Uznanie Autorstwa – Na Tych Samych Warunkach}

\title{Jak zostać mentorem w CoderDojo?}
\author{Fundacja CoderDojo Polska}
\date{\today}


\begin{document}
\maketitle
\includegraphics[width=\textwidth]{logo_CD.pdf}

\section{5 kroków, aby stać się Mentorem}

\begin{enumerate}
\item Znajdź najbliższe dojo na naszej mapie (link). Jeśli blisko Ciebie nie ma żadnego pomyśl, czy nie chciałbyś stworzyć nowego. 
\item Przyjdź do nas na zajęcia i porozmawiaj z mentorami. Spotykamy się w soboty przed południem. 
\item Pomyśl, ile czasu możesz nam poświęcić. Nie ważne, czy będziesz z nami na 100\% czy 1\%. Ważne, żebyś chciał nam pomóc!
\item Nie martw się, jeśli nie umiesz programować. Każdy się przyda. Możesz pomagać dzieciakom, możesz się od nich uczyć, możesz je wspierać i motywować. Jeśli zdecydujesz się poświęcić swój czas – znajdziemy dla Ciebie zajęcie 
\item Poproś nas o materiały. Mamy scenariusze, samouczki i dużą kolektywną wiedzę na tematy związane z robotyką i programowaniem\footnote{Część z nich znajdziesz np. na \url{https://zasoby.coderdojo.org.pl/}}. W CoderDojo uczymy się od siebie nawzajem. 
\end{enumerate}
\section{Co Ci da dołączenie do grona mentorów CoderDojo?}
\begin{itemize}
\item Będziesz miał możliwość współtworzenia najfajniejszej inicjatywy programistycznej w Polsce. 
\item Poznasz fajnych ludzi.  
\item Pomożesz dzieciakom rozwijać pasje, które mogą stać się dla nich szansą na oszałamiające kariery w stylu Marka Zuckerberga. 
\item Staniesz się częścią  społeczności CoderDojo. Będziesz się z nami spotykał na Dojo Konfach, brał udział w fajnych wydarzeniach programistycznych i robił ważne rzeczy dla przyjemności. 
\item Będziesz miał możliwość testowania sprzętów i programów, które są na stanie fundacji. 
\end{itemize}


\end{document}